\section{Limitations and Threats to Validity}\label{sec:validity}
\subsection{Selection and statistical validity}
Stage 1 screens methods on the full population, including sentences later assigned to Stage 2 test. Final testing is therefore exploratory rather than independent confirmation. Neural CV selects an epoch on the same validation fold used for its reported score. This can introduce optimistic checkpoint-selection bias and differs from the fixed classical fitting process. A future outer holdout should exclude test data from configuration and epoch selection; inner validation or a fixed training schedule would address the checkpoint issue.

The exploration budget is asymmetric: the five classical methods use fixed hyperparameter settings, with no systematic hyperparameter search documented in the archived runs, whereas the Transformer family is screened across seventeen configurations (E1--E17). Neural checkpoint selection adds another selection opportunity. The comparison therefore concerns the recorded configurations, not equally tuned model families; it does not establish that a tuned classical baseline would be inferior. A controlled follow-up should allocate comparable tuning budgets within inner validation and evaluate frozen choices on an independent outer test set.

The cohort has one base seed, 42; fold seeds 42--46 are not repeated full-protocol seeds. Screening seventeen configurations creates multiple-comparison opportunities, and fold SD is descriptive variation rather than a confidence interval. No statistical significance is claimed. Small label supports make some F1 values sensitive to a few decisions. Repeated protocol seeds and paired uncertainty analysis are needed to establish whether small differences persist.

\subsection{Data and external validity}
Exact whitespace-normalized duplicates stay in one group; fold-local vectorizers, weights and graphs avoid fitting these statistics on validation/test labels. However, near duplicates and class/project conventions remain possible. Train/test class overlap is substantial, and project identity is not a grouping variable. Generalization to new projects or developer workflows remains untested. Different language taxonomies/supports also prevent attributing corpus-level score differences to language alone. Information-type classification does not measure comment correctness or usefulness.

\subsection{Metrics and experimental coverage}
The final neural subset is E1/E3/E9/E15/E17, not all seventeen configurations or each language's saved top four. Missing final results are not inferred. E4 screening uses E1's threshold 0.5; its final threshold-tuning experiment is absent. Threshold sensitivity and score scales can affect F1 and pooled AP, motivating our precision/recall, macro-F1 and support analysis. Published scores use official splits and aggregation, as distinguished in Section~\ref{sec:published}.

\subsection{Artifact and software scope}
Aggregate/per-label metrics are available, but original sentence predictions and distributed checkpoints are not. These support numerical error counts but cannot identify actual sentence-level false positives/negatives or replay decisions. Section~\ref{sec:illustrations} labels its annotated examples as illustrations, not model errors. Exact historical package/Python versions and GPU/driver/cost measurements are missing, so runtime is not a controlled efficiency comparison.

The new adapters cover staged-dataset persistence; they do not convert every tokenizer, checkpoint or plot writer into a switchable backend. They are verified against the same data without rerunning archived experiments. Future releases should preserve predictions, model files, environment metadata and project-grouped split IDs before evaluating frozen candidates on an independent outer protocol.
